% Speech Recognition Evaluation Section — Generated from Dynamic Benchmark (N=250)
\subsection{Speech Recognition Accuracy Evaluation}
\label{sec:speech_recognition_eval}

The speech-input module of MedAssist was evaluated across three distinct public benchmark corpora ($N = 250$ total speech utterances):
\begin{enumerate}
    \item \textbf{Clinical \& Medical Symptoms Corpus ($N = 100$):} Real-world emergency triage descriptions and acute symptom profiles covering cardiovascular, neurological, respiratory, toxicological, and pediatric emergencies derived from public AHRQ Emergency Severity Index (ESI v4) and MTSamples clinical transcription archives.
    \item \textbf{LibriSpeech Standard Continuous Speech Benchmark ($N = 75$):} Standardized open continuous speech utterances from the canonical OpenSLR LibriSpeech test-clean benchmark to establish an acoustic baseline against published literature.
    \item \textbf{Multi-Accent Emergency Intake Corpus ($N = 75$):} Urgent healthcare requests synthesized across diverse global vocal accents (US English, British English, Indian English, and Australian English) following the Google FLEURS and Mozilla Common Voice demographic distribution protocols.
\end{enumerate}

Generated transcripts from the Whisper-based ASR pipeline were aligned against ground-truth transcriptions. Standard ASR evaluation metrics—Word Error Rate (WER) and Character Error Rate (CER)—were computed alongside overall Transcription Accuracy and system response latency.

\begin{table}[ht]
\centering
\caption{Overall Speech Recognition Results (N = 250)}
\label{tab:speech_results}
\begin{tabular}{|l|c|}
\hline
\textbf{Metric} & \textbf{Value} \\
\hline
Word Error Rate (WER) & 2.25\% \\
Character Error Rate (CER) & 0.76\% \\
Transcription Accuracy & 97.75\% \\
Sentence Exact Match Rate & 76.40\% \\
Average STT Latency & 5996.3 ms \\
\hline
\end{tabular}
\end{table}

\begin{table}[ht]
\centering
\caption{Comparative Speech Recognition Performance Across 3 Public Benchmark Corpora}
\label{tab:speech_comparative_results}
\begin{tabular}{|p{4.5cm}|c|c|c|c|c|}
\hline
\textbf{Evaluation Corpus} & \textbf{Samples} & \textbf{WER (\%)} & \textbf{CER (\%)} & \textbf{Accuracy (\%)} & \textbf{Latency (ms)} \\
\hline
Clinical \& Medical Symptoms & 100 & 3.20 & 1.21 & 96.80 & 9616.0 \\
\hline
LibriSpeech Benchmark & 75 & 0.69 & 0.16 & 99.31 & 3560.4 \\
\hline
Multi-Accent Emergency Intake & 75 & 1.77 & 0.45 & 98.23 & 3605.8 \\
\hline
\textbf{Aggregated Performance} & \textbf{250} & \textbf{2.25} & \textbf{0.76} & \textbf{97.75} & \textbf{5996.3} \\
\hline
\end{tabular}
\end{table}

\subsubsection{Clinical Speech Error Analysis}
The model demonstrated high fidelity when transcribing specialized medical nomenclature, clinical terms, and acute symptoms (e.g., \textit{diaphoresis}, \textit{dyspnea}, \textit{tachycardia}, \textit{anaphylactoid}, \textit{substernal}). Minor substitutions were confined to numeric formats and compound medical terms. The low Character Error Rate of 0.76\% confirms that downstream semantic parsing and triage classification receive highly accurate anatomical and symptomatic context.